% Part: computability
% Chapter: recursive-functions
% Section: pr-relations

\documentclass[../../../include/open-logic-section]{subfiles}

\begin{document}

\olfileid{cmp}{rec}{prr}
\olsection{આદિમ પુનરાવર્તી સંબંધો}


\begin{defn}
સંબંધ $R(\vec x)$નું લક્ષણ વિધેય
\[
\Char{R}(\vec x) = \left\{
  \begin{array}{ll}
  1 & \mbox{જો $R(\vec x)$ લાગુ પડે} \\
  0 & \mbox{અન્યથા}
  \end{array}
\right.
\]
આદિમ પુનરાવર્તી હોય, તો આ સંબંધને આદિમ પુનરાવર્તી
કહેવાય છે.
\end{defn}

બીજા શબ્દોમાં, આદિમ પુનરાવર્તી સંબંધ $R(\vec x)$ કહીએ ત્યારે
$\Char{R}(\vec x) = 1$ સ્વરૂપના સંબંધની વાત કરીએ છીએ. અહીં
$\Char{R}$ એવું આદિમ પુનરાવર્તી વિધેય છે જે કોઈ પણ ઇનપુટ પર
1 અથવા 0 આપે છે. દાખલા તરીકે, $x = 0$ હોય ત્યારે અને માત્ર
ત્યારે જ લાગુ પડતો સંબંધ $\fn{IsZero}(x)$, નીચેના આદિમ
પુનરાવર્તનથી વ્યાખ્યાયિત વિધેય $\Char{\fn{IsZero}}$ને અનુરૂપ છે:
\begin{align*}
\Char{\fn{IsZero}}(0) & = 1,\\
\Char{\fn{IsZero}}(x+1) & = 0.
\end{align*}

સંબંધોનું અન્ય આદિમ પુનરાવર્તી વિધેયો સાથે સંયોજન પણ કરી શકાય છે.
તેથી નીચેના સંબંધો પણ આદિમ પુનરાવર્તી છે:
\begin{enumerate}
\item સમાનતાનો સંબંધ $x = y$, જેને $\fn{IsZero}(\left|x -
  y\right|)$થી વ્યાખ્યાયિત કરીએ છીએ.
\item નાનું અથવા સમાન હોવાનો સંબંધ $x \leq y$, જેને
  $\fn{IsZero}(x \tsub y)$થી વ્યાખ્યાયિત કરીએ છીએ.
\end{enumerate}

\sourcecorrection{OLCMP-003}{સ્થિર અંગ્રેજી મૂળમાં બીજા મુદ્દાને
``less-than relation'' કહે છે, પરંતુ આપેલું સૂત્ર $x \leq y$ અને
$\fn{IsZero}(x \tsub y)$ બન્ને ``નાનું અથવા સમાન'' દર્શાવે છે.
ગુજરાતી નામમાં સૂત્ર પ્રમાણે સમાનતાનો સમાવેશ રાખ્યો છે.}

\begin{prop}
  આદિમ પુનરાવર્તી સંબંધોનો ગણ બુલિયન ક્રિયાઓ હેઠળ બંધ છે. એટલે
  જો $P(\vec x)$ અને $Q(\vec x)$ આદિમ પુનરાવર્તી હોય, તો નીચેના
  સંબંધો પણ એવા જ છે:
  \begin{enumerate}
  \item $\lnot P(\vec x)$
  \item $P(\vec x) \land Q(\vec x)$
  \item $P(\vec x) \lor Q(\vec x)$
  \item $P(\vec x) \lif Q(\vec x)$
  \end{enumerate}
\end{prop}

\begin{proof}
  ધારો કે $P(\vec x)$ અને $Q(\vec x)$ આદિમ પુનરાવર્તી છે, એટલે
  તેમનાં લક્ષણ વિધેયો $\Char{P}$ અને $\Char{Q}$ આદિમ પુનરાવર્તી છે.
  આપણે $\lnot P(\vec x)$ વગેરેના લક્ષણ વિધેયો પણ આદિમ પુનરાવર્તી છે
  તે બતાવવાનું છે.
  \[
  \Char{\lnot P}(\vec x) = \begin{cases}
    0 & \text{જો $\Char{P}(\vec x) = 1$ હોય}\\
    1 & \text{અન્યથા}
  \end{cases}
  \]
  $\Char{\lnot P}(\vec x)$ને $1 \tsub \Char{P}(\vec x)$ તરીકે
  વ્યાખ્યાયિત કરી શકીએ.
  \[
  \Char{P \land Q}(\vec x) = \begin{cases}
    1 & \text{જો $\Char{P}(\vec x) = \Char{Q}(\vec x) = 1$ હોય}\\
    0 & \text{અન્યથા}
  \end{cases}
  \]
  $\Char{P \land Q}(\vec x)$ને $\Char{P}(\vec x) \cdot
  \Char{Q}(\vec x)$ અથવા $\fn{min}(\Char{P}(\vec x), \Char{Q}(\vec
  x))$ તરીકે વ્યાખ્યાયિત કરી શકીએ. એ જ રીતે,
  \begin{align*}
    \Char{P \lor Q}(\vec x) & = \fn{max}(\Char{P}(\vec x), \Char{Q}(\vec x)) \text{ અને}\\
    \Char{P \lif Q}(\vec x) & = \fn{max}(1 \tsub
  \Char{P}(\vec x), \Char{Q}(\vec x)).
  \end{align*}
\end{proof}

\begin{prop}
  આદિમ પુનરાવર્તી સંબંધોનો ગણ સીમિત પરિમાણીકરણ હેઠળ બંધ છે. એટલે
  જો $R(\vec x, z)$ આદિમ પુનરાવર્તી સંબંધ હોય, તો નીચેના સંબંધો
  પણ એવા જ છે:
  \begin{align*}
    & \bforall{z < y}{R(\vec x, z)} \text{ અને}\\
    & \bexists{z < y}{R(\vec x, z)}.
  \end{align*}
  $\bforall{z < y}{R(\vec x, z)}$ એ $\vec x$ અને $y$ માટે ત્યારે
  અને માત્ર ત્યારે જ લાગુ પડે જ્યારે $y$થી નાની દરેક~$z$ માટે
  $R(\vec x, z)$ લાગુ પડે; $\bexists{z < y}{R(\vec x, z)}$ માટે પણ
  અનુરૂપ અર્થ છે.
\end{prop}

\begin{proof}
  રૂઢિ પ્રમાણે $\bforall{z < 0}{R(\vec x, z)}$ને સત્ય માનીએ છીએ
  (કારણ કે $0$થી નાની કોઈ $z$ છે જ નહીં) અને
  $\bexists{z < 0}{R(\vec x, z)}$ને અસત્ય માનીએ છીએ. સીમિત
  સર્વપરિમાણક સાન્ત ગુણાકાર અથવા ક્રમિક લઘુત્તમની જેમ કામ કરે છે.
  એટલે $P(\vec x, y) \defiff \bforall{z <
  y}{R(\vec x, z)}$ હોય તો $\Char{P}(\vec x, y)$ને આમ વ્યાખ્યાયિત
  કરી શકાય:
  \begin{align*}
    \Char{P}(\vec x, 0) & = 1\\
    \Char{P}(\vec x, y+1) & =
    \fn{min}(\Char{P}(\vec x, y), \Char{R}(\vec x, y)).
  \end{align*}
  સીમિત અસ્તિત્વપરિમાણીકરણને એ જ રીતે $\fn{max}$ વડે વ્યાખ્યાયિત
  કરી શકાય. અથવા તેને
  $\bexists{z < y}{R(\vec x, z)}
  \liff \lnot \bforall{z < y}{\lnot R(\vec x, z)}$ સમકક્ષતા
  વાપરીને સીમિત સર્વપરિમાણીકરણમાંથી મેળવી શકાય. નોંધો કે,
  દાખલા તરીકે, $\bexists{x \leq y}{\dots
  x\dots}$ સ્વરૂપનો સીમિત પરિમાણક
  $\bexists{x < y+1}{\dots x \dots}$ને સમકક્ષ છે.
\end{proof}

\begin{prob}
  ત્રણ આર્ગ્યુમેન્ટવાળો સંબંધ $x \equiv y \mod n$ (મોડ્યુલો~$n$
  મુજબ સર્વાંગસમતા) આદિમ પુનરાવર્તી છે એવું બતાવો.
\end{prob}

બીજું ઉપયોગી આદિમ પુનરાવર્તી વિધેય શરતી વિધેય
$\fn{cond}(x,y,z)$ છે, જે આમ વ્યાખ્યાયિત છે:
\begin{align*}
  \fn{cond}(x,y,z) & = \begin{cases}
  y & \text{જો $x = 0$ હોય} \\
  z & \text{અન્યથા}.
\end{cases}
\intertext{આ વિધેયને પુનરાવર્તી રીતે આમ વ્યાખ્યાયિત કરીએ:}
\fn{cond}(0,y,z) & = y,\\
\fn{cond}(x+1,y,z) & = z.
\end{align*}
આનો ઉપયોગ કરીને આદિમ પુનરાવર્તી સંબંધોના આધારે કિસ્સાવાર
વ્યાખ્યાયિત વિધેયોનું સમર્થન કરી શકીએ:

\begin{prop}
જો $g_0(\vec x)$, \dots,~$g_m(\vec x)$ આદિમ પુનરાવર્તી વિધેયો
હોય અને $R_0(\vec
x)$, \dots, $R_{m-1}(\vec x)$ આદિમ પુનરાવર્તી સંબંધો હોય, તો
નીચે મુજબ વ્યાખ્યાયિત વિધેય $f$ પણ આદિમ પુનરાવર્તી છે:
\[
f(\vec x) = \begin{cases}
    g_0(\vec x) & \text{જો $R_0(\vec{x})$ લાગુ પડે} \\
    g_1(\vec x) & \text{જો $R_1(\vec{x})$ લાગુ પડે અને $R_0(\vec{x})$ ન પડે} \\
    \vdots & \\
    g_{m-1}(\vec x) & \text{જો $R_{m-1}(\vec{x})$ લાગુ પડે અને અગાઉના
      કોઈ ન પડે}
    \\
    g_m(\vec x) & \mbox{અન્યથા}
\end{cases}
\]
\end{prop}

\begin{proof}
  $m = 1$ હોય ત્યારે આ માત્ર નીચે મુજબ વ્યાખ્યાયિત વિધેય છે:
  \[
  f(\vec x) = \fn{cond}(\Char{\lnot R_0}(\vec x),g_0(\vec x),g_1(\vec
  x)).
  \]
  $m$ એ $1$થી મોટું હોય ત્યારે આ સ્વરૂપની વ્યાખ્યાઓનું સંયોજન
  કરી શકાય છે.
\end{proof}
\end{document}
